\section*{7. Autoevaluación Docente}

Mi enfoque docente se basa en tres pilares fundamentales: intergeneracional, internacional e interdisciplinario, los cuales enriquecen tanto el aprendizaje como la investigación:

\begin{itemize}
    \item \textbf{Intergeneracional:} Colaboro con mentores senior, quienes guían mi desarrollo académico, mientras que al mismo tiempo actúo como mentor para estudiantes de diversos niveles, incluyendo tesis de doctorado, maestría, pregrado e incluso proyectos escolares. Este ciclo de aprendizaje continuo me permite integrar experiencia y energía en la formación académica.
    
    \item \textbf{Internacional e Intercultural:} Mi experiencia en instituciones de primer nivel en Chile, Suiza y los Países Bajos ha consolidado mi capacidad para colaborar con estudiantes y académicos de diversos contextos culturales. Estas interacciones fomentan un ambiente académico inclusivo y enriquecen las perspectivas de mis estudiantes.
    
    \item \textbf{Interdisciplinario:} La versatilidad de las matemáticas como lenguaje universal me ha permitido explorar múltiples campos, desde biología y ciencias ambientales hasta ingeniería y finanzas. Este enfoque interdisciplinario enriquece mi enseñanza al conectar la teoría matemática con aplicaciones concretas y variadas.
\end{itemize}

Mi compromiso con la docencia se manifiesta en la creación de materiales de aprendizaje rigurosos y accesibles, como apuntes detallados y grabaciones de clase. También superviso tesis y proyectos que integran teoría y práctica, fomentando la curiosidad y el análisis crítico en los estudiantes.

Como educador, mi objetivo es inspirar a los estudiantes a trabajar con rigor y aplicar su conocimiento a problemas reales, desarrollando habilidades críticas y creativas. A través de proyectos interdisciplinarios y colaboraciones internacionales, los preparo para ser agentes de cambio tanto en sus comunidades como en el ámbito científico global.